% !TeX spellcheck = es_ES
\documentclass[letterpaper,11pt]{article}
\usepackage[utf8]{inputenc}
\usepackage[T1]{fontenc}
\usepackage[spanish]{babel}
\decimalpoint % separador decimal con punto, como en las tablas de la fuente original
\usepackage{lmodern,microtype}
\usepackage{amsmath,amssymb,amsfonts}
\usepackage{geometry}
\geometry{margin=2.54cm}
\usepackage{setspace}
\setstretch{1.15}
\usepackage{enumitem}
\usepackage{booktabs}
\usepackage{color}
\usepackage{hyperref}
\definecolor{ntr}{rgb}{0.8,0,0}
\hypersetup{colorlinks=true,linkcolor=ntr,citecolor=ntr,urlcolor=black}

% Portada (Luis Cabral): titulo \LARGE sans-serif alineado a la izquierda, sin numero de pagina
\makeatletter
    \renewcommand{\@maketitle}{%
    \newpage\parindent0in\null\vskip2em%
    {\LARGE\textbf{\textsf{\@title}}\par\vskip2em}%
    \@author\par\vskip2em%
    \@date%
    \par
    \thispagestyle{empty}\setcounter{page}{1}%
    }%
\makeatother

\DeclareMathOperator{\Var}{Var}
\DeclareMathOperator{\Cov}{Cov}
\DeclareMathOperator{\Corr}{Corr}
\DeclareMathOperator{\E}{\mathbb{E}}
\DeclareMathOperator{\N}{Normal}
\DeclareMathOperator{\LP}{\mathbb{L}}
\newcommand\independent{\protect\mathpalette{\protect\independenT}{\perp}}
\def\independenT#1#2{\mathrel{\rlap{$#1#2$}\mkern2mu{#1#2}}}
% Teoria asintotica
\newcommand{\SR}{\mathbb{R}}
\DeclareMathOperator{\plim}{plim}
\DeclareMathOperator{\Avar}{AVar}
\newcommand{\pto}{\stackrel{p}{\longrightarrow}}
\newcommand{\dto}{\stackrel{d}{\longrightarrow}}
\newcommand{\adistr}{\stackrel{a}{\sim}}
\newcommand{\epsi}{\varepsilon}


\title{Análisis Econométrico \\ Ayudantía 8}
\author{Magíster en Economía, Universidad Alberto Hurtado}
\date{}

% Funciones de R sugeridas: entre corchetes, alineadas a la derecha; si no caben en la
% linea pasan enteras a la siguiente.
\newcommand{\fn}[1]{{\unskip\nobreak\hfill\penalty50\hskip1em\null\nobreak\hfill
    \mbox{[#1]}\parfillskip=0pt\finalhyphendemerits=0\par}}

\begin{document}
\maketitle

\noindent Temas: cómo hacer inferencia en R (errores estándar robustos, estadístico t,
p-valor e intervalos de confianza), contrastes e intervalos para una combinación lineal de
coeficientes, intervalos para una transformación no lineal (método delta) y un
primer vistazo al test de Wald con dos restricciones.

\medskip
\noindent Los ejercicios 1 y 2 son una demostración en R: el ayudante corre el código en
pantalla, así que no hace falta traer computador. El ejercicio 3 se resuelve en la pizarra si
queda tiempo; si no, queda como práctica. La parte 2(d) adelanta el test de Wald, que se ve en la
clase del lunes 2 de noviembre.

\section*{Si quiere replicarlo en su computador}

\noindent El script completo, \texttt{TA8.R}, y los datos están en el repositorio del curso,
en \href{https://github.com/ramirode/econometria_mae/tree/main/TA/TA8}{\texttt{TA/TA8}}. Los
pasos son los mismos de la Ayudantía 5: cree una carpeta \texttt{TA8} con las subcarpetas
\texttt{code} (para \texttt{TA8.R}) e \texttt{input} (para \texttt{caschool.dta} e
\texttt{injury.dta}), y cambie la
línea \texttt{dir\_proyecto <- "..."} por la ruta a su carpeta \texttt{TA8}. En Posit Cloud la
carpeta ya existe: \texttt{/cloud/project/TA/TA8}. Paquetes: \texttt{haven},
\texttt{sandwich}, \texttt{lmtest} y \texttt{car}.

\section*{Ejercicios}

\begin{enumerate}[itemsep=0.8em]

%% ---------------------------------------------------------------- 1
\item \emph{Inferencia en R: tamaño de clase y rendimiento académico.} Con los datos
    \texttt{caschool.dta} (420 distritos escolares de California, 1998--1999), estime
    \begin{align*}
    testscr_i=\beta_0+\beta_1\, str_i+\beta_2\, el\_pct_i+u_i,
    \end{align*}
    donde $testscr$ es el puntaje promedio en una prueba estandarizada, $str$ es el número de
    alumnos por profesor y $el\_pct$ es el porcentaje de alumnos que están aprendiendo inglés.
    Son los datos del ejemplo de la clase: compare sus resultados con las diapositivas.
    \begin{enumerate}
        \item Estime el modelo y mire la salida de \texttt{summary}. ¿Qué errores estándar
              reporta? Calcule los errores estándar robustos a heteroscedasticidad y compárelos.
              \fn{\texttt{lm}, \texttt{summary}, \texttt{vcovHC}, \texttt{coeftest}}
        \item Contraste $H_0:\beta_1=0$ contra la alternativa bilateral. Calcule el estadístico
              t y el p-valor a mano, con el coeficiente y la matriz de varianzas y covarianzas, y
              compárelos con la salida de \texttt{coeftest}.
              \fn{\texttt{coef}, \texttt{sqrt}, \texttt{pnorm}}
        \item Construya un intervalo de confianza al 95\% para $\beta_1$, a mano y con
              \texttt{coefci}.
              \fn{\texttt{qnorm}, \texttt{coefci}}
        \item Contraste $H_0:\beta_1=0$ contra la alternativa de que el tamaño de clase tiene un
              efecto negativo, $H_1:\beta_1<0$. Calcule el valor crítico y el p-valor.
              \fn{\texttt{qnorm}, \texttt{pnorm}}
    \end{enumerate}

%% ---------------------------------------------------------------- 2
\item \emph{Compensación por accidentes laborales (R).} Meyer, Viscusi y Durbin
    (\textit{American Economic Review}, 1995) estudian el efecto de la compensación semanal que
    reciben los trabajadores accidentados sobre la duración de la compensación (en semanas). La
    compensación es proporcional al salario, hasta un tope. En Kentucky, el 15 de julio de 1980
    el tope subió de 131 a 217 dólares por semana: la compensación aumentó para los
    trabajadores de altos ingresos y no cambió para los de bajos ingresos.
    Con los datos \texttt{injury.dta}, sólo para Kentucky ($ky=1$), considere
    \begin{align*}
    \log(durat_{i})=\beta_0+\beta_1\, afchnge_i+\beta_2\, highearn_{i} +\beta_3\, afhigh_{i}+u_i,
    \end{align*}
    donde $afchnge_i$ es igual a 1 si el accidente ocurrió después del cambio,
    $highearn_i$ es igual a 1 si el trabajador es de altos ingresos y
    $afhigh_i=afchnge_i \times highearn_i$.
    \begin{enumerate}
        \item Estime el modelo con errores estándar robustos. Escriba la media de
              $\log(durat)$ en los cuatro grupos en función de los $\beta$'s. ¿Qué mide
              $\beta_3$?
              \fn{\texttt{lm}, \texttt{vcovHC}, \texttt{coeftest}}
        \item Para los trabajadores de altos ingresos, estime la diferencia en la media de
              $\log(durat)$ después y antes del cambio; llámela $\theta$. Contraste
              si es distinta de cero de dos formas: a mano, con la matriz de varianzas y
              covarianzas, y con \texttt{linearHypothesis}.
              \fn{\texttt{vcovHC}, \texttt{linearHypothesis}}
        \item Exprese la diferencia de (b) como cambio porcentual exacto en la duración,
              $100\,(\exp(\theta)-1)$, y construya un intervalo de confianza al 95\% con el
              método delta, a mano y con \texttt{deltaMethod}.
              \fn{\texttt{exp}, \texttt{deltaMethod}}
        \item \emph{Adelanto: test de Wald.} Contraste la hipótesis nula de que el cambio del
              tope no afectó la duración \emph{ni para los de altos ingresos ni para los de
              bajos ingresos}. ¿Cuántas restricciones hay? Escriba $R'$ y calcule el
              estadístico de Wald a mano y con \texttt{linearHypothesis}.
              \fn{\texttt{solve}, \texttt{pchisq}}
    \end{enumerate}

%% ---------------------------------------------------------------- 3
\item \emph{Trabajo en equipo y riesgo moral (Examen 2016).} David C. Chan, en
    \textit{``Teamwork and Moral Hazard: Evidence from the Emergency Department''} (Journal
    of Political Economy, 2016), estudia cómo la organización del trabajo afecta la
    productividad, con datos del servicio de urgencias de un hospital con dos salas. En una, las
    enfermeras asignan cada paciente a un médico determinado. En la otra, el sistema \emph{de
    trabajo en equipo}, las enfermeras asignan el paciente a una cama que es responsabilidad de
    dos médicos, quienes deciden entre ellos quién lo atiende. La hipótesis es que, cuando las
    enfermeras asignan los pacientes, los médicos tienen incentivos a \emph{arrastrar los pies}
    (demorarse más de lo necesario con cada paciente para recibir menos pacientes nuevos), y que
    el trabajo en equipo mitiga esos incentivos.

    Suponga que los médicos no eligen en qué sala trabajan y que no hay otras amenazas a la
    validez interna. El modelo es
    \begin{align*}
    \log(length\_of\_stay_{i})=\beta_0+\beta_1 \,self\_managed_{i}+x_{i}'\gamma+u_{i},
    \end{align*}
    donde $i$ es una atención de urgencia, $length\_of\_stay_i$ es el tiempo de atención (en
    horas), $self\_managed_{i}$ es igual a 1 si el paciente se atiende en la sala con trabajo
    en equipo, y $x_{i}$ incluye características del paciente (edad, sexo, seguro de salud,
    severidad de la urgencia, diagnóstico, etc.).
    \begin{enumerate}
        \item La estimación es $\hat \beta_1 = -0.147$, con $se(\hat \beta_1) = 0.046$.
              Interprete el resultado. Contraste si el efecto es estadísticamente significativo
              y construya un intervalo de confianza al 95\%. ¿Es económicamente significativo?
        \item Para ver si los incentivos a arrastrar los pies son mayores cuando hay más
              trabajo esperado, el autor estima
              \begin{align*}
              \log(length\_of\_stay_{i})=&\beta_0+\beta_1 \, self\_managed_{i}+\beta_2 \, waiting_{i}\\
              &+\beta_3 \, self\_managed_{i} \times waiting_{i} +x_{i}'\gamma+u_{i},
              \end{align*}
              donde $waiting_i$ es el número de pacientes en la sala de espera. Los resultados
              son $\hat \beta_1 = -0.0101$, $se(\hat \beta_1) = 0.0196$,
              $\hat \beta_2 = 0.0063$, $se(\hat \beta_2) = 0.0013$, $\hat \beta_3 =-0.0060$ y
              $se(\hat \beta_3) = 0.0014$. Contraste si el efecto del trabajo en equipo depende
              del número de pacientes en la sala de espera.
        \item Escriba el efecto del trabajo en equipo para un valor dado $w$ de $waiting$,
              $\theta(w)$, en función de los $\beta$'s. Proponga un intervalo de confianza para
              $\theta(w)$ con nivel de confianza asintótico del 95\%, expresado en función de
              $w$, los $\hat\beta$'s y los elementos de la matriz de varianzas y covarianzas
              estimada. Demuestre que su nivel de confianza asintótico es 95\%, enunciando los
              teoremas que utiliza.
        \item Con $\widehat{\Avar}(\hat \beta_1,\hat \beta_3)=-0.000014$, calcule el intervalo
              de (c) cuando no hay pacientes en la sala de espera ($w=0$) y cuando hay 14.1
              (la media más una desviación estándar). Interprete.
    \end{enumerate}
\end{enumerate}

\end{document}
